\tikzset{add/.style n args={4}{
    minimum width=14,
    path picture={
    \draw[black] 
        (path picture bounding box.south east) -- (path picture bounding box.north west)
        (path picture bounding box.south west) -- (path picture bounding box.north east);
    \node[anchor=north, inner sep=0pt] at ($(path picture bounding box.north)+(0,-0.13)$) {\tiny #1};
    \node[anchor=east,  inner sep=0pt] at ($(path picture bounding box.east)+(-0.13,0)$)  {\tiny #2};
    \node[anchor=south, inner sep=0pt] at ($(path picture bounding box.south)+(0,0.13)$) {\tiny #3};
    \node[anchor=west,  inner sep=0pt] at ($(path picture bounding box.west)+(0.13,0)$)  {\tiny #4};
        }
    }
}

\tikzstyle{block} = [draw, rectangle, align=center, font=\small]%
\tikzstyle{sum} = [draw, circle, add={}{$+$}{$+$}{$-$}]%
\tikzstyle{input} = [text centered, inner sep=0pt, outer sep=0pt, font=\footnotesize]%
\tikzstyle{output} = [text centered, inner sep=0pt, outer sep=0pt, font=\footnotesize]%

\begin{tikzpicture}[auto,node distance=1cm,>=latex']
	% ----------------------------------
	% --- définition des espacements ---
	% ----------------------------------
	\def\x{2.5em}
	\def\y{2em}
    \def\dx{-250}
    \def\dy{1}

    % ----------------------------
	% --- définition des nodes ---
	% ----------------------------
	% node de consigne
	\node [input] (r2) {$r_2(k)$};

    % node des régulateurs + préfiltre
    \node [block, below=\y of r2.south, fill=cyan!30, anchor=north] (CpCiPf) {$C_P(q)+C_I(q)P\!F(q)$}; 

    % node sommateur
    \node [sum, below=\y of CpCiPf.south, fill=cyan!30, anchor=north] (sum3) {};

    % node des régulateurs
    \node [block, right=\x of sum3.east, fill=cyan!30, anchor=west] (CpCi) {$C_P(q)+C_I(q)$}; 

    % node de commande
    \node [input] (r1) at ([xshift=-\x]r2-|CpCiPf.west) {$r_1(k)$};

    % node de sortie
    \node [input] (y) at ([xshift=\x]r2-|CpCi.east) {$y(k)$};

    % node filtres
    \node [block, below=\y of sum3.south, fill=red!30, anchor=north] (Fsvfl) {$F_{svf}(q)$};
    \node [block, fill=red!30] (Fsvfr) at (Fsvfl-|y) {$F_{svf}(q)$};

    % node vecteur régresseur
    \node [block, fill=cyan!30, anchor=north] (phi) at ([yshift=-\y]$(Fsvfl.south)!0.5!(Fsvfr.south)$){%
        $\begin{aligned}%
              \varphi_f^T(k) = [ - & y_f(k-n_a),\,\cdots, -y_f(k-1),                       \\%
                             & u_f(k-n_b-n_k-1), \,\cdots, \,u_f(k-n_k) ]%
          \end{aligned}$%
    };

    % node matrice de régression
    \node [block, below=\y of phi.south, fill=cyan!30, anchor=north] (Phi) {%
        $\Phi_f = %
        \begin{bmatrix} %
            \varphi_f^T(1) \\ %
            \vdots \\ %
            \varphi_f^T(N) %
        \end{bmatrix}$%
    };

    % node matrice de sortie
    \node [block, below right=\y and \x of phi.south east, fill=cyan!30, anchor=north] (Y) {%
        $Y_f = %
        \begin{bmatrix} %
            y_f(1) \\ %
            \vdots \\ %
            y_f(N) %
        \end{bmatrix}$%
    };

    % node solution
    \node [block, fill=cyan!30, anchor=north] (sol) at ([yshift=-\y]$(Phi.south)!0.5!(Y.south)$){\textbf{Solution CLLSSVFFR}\\$\hat{\theta}_{ls}=\left(\Phi_f^T \Phi_f\right)^{-1}\Phi_f^T Y_f$};

    % node de sortie
    \node [output, below=\y of sol.south, anchor=north] (theta) {$\hat{\theta}_{ls}$};

    % ------------------------------
	% --- définition des flèches ---
	% ------------------------------
    % de r1 au sommateur
    \draw [->] (r1) |- (sum3);

    % de r2 au sommateur
    \draw [->] (r2) -- (CpCiPf);
    \draw [->] (CpCiPf) -- (sum3);

    % de y au sommateur
    \draw [->] (y) |- (CpCi);
    \draw [->] (CpCi) -- (sum3);

    % du sommateur au filtre 
    \draw [->] (sum3) -- (Fsvfl) node [midway, right] {$u(k)$};

    % de la sortie y au filtre
    \draw [->] (CpCi-|y) -- (Fsvfr);

    % des filtres à phi
    \draw [->] (Fsvfl) -- (phi.north-|Fsvfl) node [midway, right] {$u_f(k)$};
    \draw [->] (Fsvfr) -- (phi.north-|Fsvfr) node [midway, left, name=yf] {$y_f(k)$};

    % de phi à Phi
    \draw [->] (phi) -- (Phi);

    % de yf à Y
    \draw [->] (yf) -| (Y);

    % de Phi à la solution
    \draw [->] (Phi) |- (sol);

    % de Y à la solution
    \draw [->] (Y) |- (sol);

    % de sol à theta
    \draw [->] (sol) -- (theta);

    % partie génération du signal
    \coordinate (p1) at ([xshift=\dx, yshift=-\dy]CpCi.south);
    \coordinate (p2) at ([yshift=-\dy]CpCi.south-|Y.east);
    \draw [dashed] (p1) -- (p2);
    \node [align=center, anchor=west] (sig_gen) at ($(r1.north-|p1)!0.5!(p1)$) {Génération\\du\\signal};

    % partie filtrage
    \coordinate (p3) at ([yshift=\dy]p1|-phi.north);
    \coordinate (p4) at ([yshift=\dy]p2|-phi.north);
    \draw [dashed] (p3) -- (p4);
    \coordinate (ptextfilt) at ($(p3)!0.5!(p1)$);
    \node (filt) at (ptextfilt-|sig_gen) {Filtrage};

    % partie régresseur
    \coordinate (p5) at ([yshift=-\dy]p1|-Phi.south);
    \coordinate (p6) at ([yshift=-\dy]p2|-Phi.south);
    \draw [dashed] (p5) -- (p6);
    \coordinate (ptextreg) at ($(p5)!0.5!(p3)$);
    \node [align=center] (reg) at (ptextreg-|sig_gen) {Conception\\du\\regresseur};

    % partie résolution
    \coordinate (ptextresol) at ($(p5)!0.5!(theta)$);
    \node [align=center] (resol) at (ptextresol-|sig_gen) {Résolution};
\end{tikzpicture}